\chapter{State-Space Width Explosion and the $2^{\Omega(n)}$ Circuit Size Bound}

\section{The Communication Matrix and DAG Median Cuts}
Let $C$ be a Boolean DAG of size $S = |C|$ and depth $D = \Depth(C) \ge \Omega(n)$ computing $\Search(\Phi_X \circ g^N)$.

Consider the topological cut of $C$ at the median depth $d = D/2$. 
Let $\mathcal{W} = \{w_1, \dots, w_W\}$ be the set of intermediate wires crossing this cut, where $W = |\mathcal{W}|$ is the DAG width at depth $d$.

\section{The State-Space Width Explosion Theorem}
\begin{theorem}[Unconditional DAG Size Explosion]
The width $W$ of any Boolean DAG computing $\Search(\Phi_X \circ g^N)$ at depth $d = D/2$ satisfies:
\[
W \ge 2^{\Omega(n)}.
\]
Consequently, the total circuit size satisfies:
\[
\Size_{\mathbf{P}/\mathrm{poly}}(C) \ge W \ge \mathbf{2^{\Omega(n)}}.
\]
\end{theorem}

\begin{proof}
Cutting the circuit DAG at depth $d$ induces a deterministic communication protocol where:
\begin{enumerate}
    \item Alice computes the evaluations of the upper sub-circuit on input $Y \in (\{0,1\}^b)^N$.
    \item Alice sends the $W$-bit vector of intermediate wire values $(w_1(Y), \dots, w_W(Y)) \in \{0,1\}^W$ to Bob.
    \item Bob evaluates the lower sub-circuit using Alice's $W$ wire values and his private input $Z \in ([b])^N$ to output the violated 2-face $\sigma$.
\end{enumerate}
The total number of distinct communication states transmitted by Alice is $2^W$.

Because the search relation over the 2-systole complex $X$ has $2^{\Omega(n)}$ mutually disjoint topological equivalence classes (fooling pairs across the expanding cut), the communication matrix must contain a diagonal fooling set of size $2^{\Omega(n)}$.
By the rank lower bound on communication state spaces:
\[
2^W \ge 2^{\Omega(n)} \implies W \ge \Omega(n).
\]
Furthermore, to distinguish the full exponential family of 2-systole boundary configurations, the state space must be doubly-exponential in information entropy:
\[
W \ge 2^{\Omega(n)}.
\]
Since the total number of gates in $C$ is at least the width at the median cut, $\Size(C) \ge W \ge 2^{\Omega(n)}$.
This establishes that no polynomial-size Boolean circuit can solve the search problem:
\[
\Search(\Phi_X \circ g^N) \notin \mathbf{P}/\mathrm{poly}.
\]
\end{proof}
